% ============================================================
%  CONCLUSION (end of chapter)                    (Member 1)
% ============================================================
\section*{Conclusion}
\addcontentsline{toc}{section}{Conclusion}

Now look again at Meera's fees. In 2026, dearer oil raised India's demand for dollars, while foreign investors selling Indian assets cut the supply \seesec{sec:moves}. The RBI sold dollars from its reserves to slow the rupee's fall but did not stop it, and the bill for her fees went up by ₹2.5 lakh. The same move gave Ravi's family more rupees for every dirham. An exchange rate is just a price, but so much of India's trade, savings and borrowing passes through it that this one price reaches almost every household in the country.

% [added at merge] summary box answering the five questions from the Introduction
\begin{keypoints}[unbreakable]{What this chapter has shown}
\begin{itemize}
  \item \textbf{Why a currency has a price.} When goods, services or savings cross a border, one currency must be swapped for another. The exchange rate is the price of that swap (Section~\ref{sec:exrate}).
  \item \textbf{Who buys and sells currencies.} Importers, students abroad and Indians investing abroad demand dollars; exporters, workers abroad, foreign tourists and foreign investors supply them, through banks regulated by the RBI (Sections~\ref{sec:forex} and~\ref{sec:bop}).
  \item \textbf{What sets the price, and why it keeps changing.} The demand for and supply of dollars. Inflation, interest rates, capital flows, oil prices and expectations keep shifting the two curves (Sections~\ref{sec:ds}, \ref{sec:real} and~\ref{sec:moves}).
  \item \textbf{How it is managed.} Countries choose fixed, floating or managed rates. Since 1993 India has used a managed float, with the RBI using its reserves to smooth sharp swings, though no country's reserves can hold a rate for ever (Sections~\ref{sec:systems}, \ref{sec:reserves} and~\ref{sec:crises}).
  \item \textbf{Who gains and who loses.} A weaker rupee helps exporters and families receiving remittances, but hurts importers, students abroad and anyone with dollar debts, and it can push up prices at home (Sections~\ref{sec:appdep} and~\ref{sec:matters}).
\end{itemize}
\textbf{Where next?} The balance of payments, monetary policy and international trade each build on the ideas in this chapter.
\end{keypoints}
